\chapter*{Wstęp}
\addcontentsline{toc}{chapter}{Wstęp}
\markboth{Wstęp}{Wstęp}

Ta książka ma jeden cel: doprowadzić Cię, krok po kroku i w ciągu roku, od zera do poziomu
\textbf{B2} niemieckiego, na którym da się pracować jako inżynier oprogramowania w Zurychu
--- rozmawiać z rekruterem i managerem, uczestniczyć w spotkaniach technicznych, pisać e-maile
i dokumentację, a po pracy załatwić mieszkanie, ubezpieczenie i urząd.

\section*{Czego ta książka nie obiecuje}

\begin{uwaga}[Uczciwa arytmetyka]
Orientacyjne nakłady pracy dorosłego, który zaczyna od zera, to (łącznie, od początku nauki):
A1 --- 100--150 h, A2 --- 200--300 h, B1 --- 350--500 h, B2 --- 550--800 h, C1 --- 900--1200 h.
To \emph{przybliżone} wartości zebrane w rozmowie planistycznej, nie wynik badania, a rzeczywisty
czas zależy od intensywności i od kontaktu z językiem.
\end{uwaga}

Wnioski, z którymi ta książka się nie kłóci:

\begin{itemize}
\item \textbf{Godzina dziennie to około 365 godzin w roku.} To wystarcza na solidne B1, nie na B2.
  Dojście do B2 w 12 miesięcy wymaga około \textbf{600 godzin}, czyli ok.\ 100 minut dziennie
  w tygodniu i 3 godzin w weekend. Książka jest zaprojektowana na ten wariant (\emph{tryb 100}).
\item Jeśli możesz poświęcić tylko 60 minut dziennie, używaj książki w \emph{trybie 60}: te same
  jednostki, ale jedna na 1,5--2 tygodnie. Spodziewaj się B1 po roku i B2 po 18--20 miesiącach.
\item Książka nie zastąpi rozmowy z żywym człowiekiem, słuchania i czytania autentycznych tekstów
  ani egzaminu. Każdy tydzień ma zadanie \emph{produkcyjne} (mówienie albo pisanie), bo
  rozpoznawanie słów nie daje płynności.
\item \textbf{Treść nie została sprawdzona przez native speakera.} Piszę po niemiecku poprawnie
  według mojej najlepszej wiedzy, ale nikt tego jeszcze nie zweryfikował, a szwajcarskie realia
  (prawo, urzędy) zmieniają się. Fakty, które wymagają źródła pierwotnego, są oznaczone
  \verify{tak jak ten} i zliczone w Dodatku~F. Przed poważną decyzją sprawdź je u źródła.
\end{itemize}

\section*{Dlaczego Hochdeutsch najpierw}

W biurze, w dokumentach, na egzaminach i w mediach używa się standardowego niemieckiego
(\emph{Hochdeutsch}). W Szwajcarii mówi się na co dzień po szwajcarsku niemiecku
(\emph{Schweizerdeutsch}), który jest dialektem, a nie wersją Hochdeutsch z akcentem. Zasada tej
książki: najpierw standard (A1--B2), dialekt jako \emph{rozumienie}, nie jako mówienie.
Szwajcarzy zwykle przechodzą na standard, gdy rozmawiają z kimś, kto go uczy się jako obcego
języka. Dodatek~D zbiera różnice (m.in.\ \de{ss} zamiast \de{ß}, \de{Velo}, \de{Grüezi}).

\section*{Jak zbudowana jest książka}

\begin{itemize}
\item \textbf{\val{units.total} jednostek = \val{units.total} tygodni}, w czterech częściach (A1, A2, B1, B2).
  Po każdym poziomie jest tydzień buforowy: powtórka, test, odrobienie zaległości. Razem
  \val{weeks.total} tygodni.
\item Każda jednostka ma stały układ: cele \emph{can-do}, słownictwo, gramatyka z przykładami,
  dialogi, ćwiczenia z kluczem, zadanie weekendowe.
\item \textbf{Liczby są liczone, nie wpisywane.} Liczba słów, ćwiczeń i dialogów wylicza skrypt
  z danych (\code{tools/gen.py}); tabele słownictwa i pliki Anki są generowane z tych samych
  plików, więc nie mogą się rozjechać.
\end{itemize}

\begin{uwaga}[Kolory rodzajników]
\artm{der} (rodzaj męski), \artf{die} (żeński), \artn{das} (nijaki). Rodzajnik jest zawsze
wypisany --- kolor tylko pomaga zapamiętać. Rzeczownika nie uczysz się bez rodzajnika.
\end{uwaga}

\section*{Skąd się biorą materiały zewnętrzne}

Książka celowo jest szkieletem na rok, nie całą biblioteką. Do kursu \emph{DW Nicos Weg}
(A1--B1, darmowy), do fiszek \emph{Anki} i do rozmów ze mną (AI) albo z korepetytorem dochodzą
materiały z Dodatku~E. Dla B2 oficjalne listy słów nie istnieją w takiej postaci jak dla A1--B1,
więc słownictwo B2 trzeba budować z autentycznych tekstów --- stąd tak duży nacisk na czytanie
i słuchanie w częściach B1--B2.

\section*{Dla kogo}

Dla osoby dorosłej, znającej angielski i pracującej z kodem. Przykłady zawodowe pochodzą ze
świata .NET, Azure, architektury i zespołów zwinnych. Jeśli jesteś kimś innym, gramatyka i
słownictwo codzienne pozostają takie same; zmienią się tylko ramki \emph{W pracy}.
